\section{Team information}

\subsection{Members}
\label{subsec:members}

Team \textbf{RushDL - 13}, class \textbf{24KTPM2}, has two members:

\begin{center}
\captionof{table}{Team members}
\label{tab:members}
\begin{tabularx}{0.9\textwidth}{|C{1cm}|Y|C{3.3cm}|}
  \hline
  \textbf{No.} & \textbf{Full name} & \textbf{Student ID}\\
  \hline
  1 & Nguyễn Sơn Hải & 24127163\\
  \hline
  2 & Châu Vũ Trung & 24127256\\
  \hline
\end{tabularx}
\end{center}

\subsection{Team goals}

The team builds the Chat System project to apply what the Java Application
Programming course teaches to a well-organised application that covers user
management and messaging. Through the project, the members practise analysing
requirements, designing an application and developing software together.

In Version 1 (20\%), the team defines the scope of the system, designs the user
interface of both subsystems (administrator and user) and designs a database that
fits the functions in the assignment. These results are the foundation for
building and completing the functions in the next version. As this stage
requires, the screens do not need event handling yet.

For the way of working, the team aims at a clear split of responsibilities, work
finished on time and content that fully meets the assignment. The team uses
Messenger to discuss and report progress; Google Drive to store and share
documents and to finish the report together; and Git and GitHub to manage the
source code, track changes and record each member's contribution.

\noindent Where the team keeps its work:
\begin{itemize}
  \item \textbf{Documents and report:}
    \href{https://drive.google.com/drive/folders/1Kci4Ug_HXLRbBBFc0h6lVz0RbuvbcJTF?usp=sharing}
    {the team's Google Drive folder}.
  \item \textbf{Source code:}
    \href{https://github.com/nguyenshai/Project-java}
    {the team's GitHub repository}.
\end{itemize}

\clearpage
\subsection{Member roles and responsibilities}
\label{subsec:member-roles}

The team splits responsibilities by subsystem. Châu Vũ Trung focuses on the
administration and monitoring functions; Nguyễn Sơn Hải focuses on the functions
for users and their communication. This split makes each member's scope clear
and limits overlapping work.

\begin{center}
\small
\captionof{table}{Scope of each member}
\label{tab:member-roles}
\begin{tabularx}{\textwidth}{|L{3cm}|L{2.8cm}|Y|}
  \hline
  \textbf{Member} & \textbf{Scope} &
  \textbf{Main responsibilities in V1}\\
  \hline
  Nguyễn Sơn Hải\newline 24127163 &
  User subsystem &
  Designs the User subsystem interface; prepares its report content
  and screenshots.\\
  \hline
  Châu Vũ Trung\newline 24127256 &
  Administrator subsystem &
  Designs the Admin subsystem interface; prepares its report content
  and screenshots; assembles the V1 report.\\
  \hline
\end{tabularx}
\end{center}

Both members design the database together, agree on how the project is organised
and cross-check each other's work. Each member is responsible for accurate content
in their own part; the person who assembles the report keeps the presentation
consistent and checks that every section is complete.

This assignment is the initial plan and can change with the actual workload once
both members agree. The tasks and related functions are detailed in
Section~\ref{subsec:task-plan}.

\clearpage
\subsection{Working rules}

The team follows these rules:

\begin{enumerate}
  \item \textbf{Follow the requirements:}
    Before starting, each member makes clear what must be done and delivered.
    The work must meet the matching requirements of the assignment, without
    leaving anything out or changing the scope on their own.

  \item \textbf{Meet deadlines:}
    Every task has a deadline the team agrees on, which fits the course
    deadline and leaves time for review and fixes.

  \item \textbf{Report progress early:}
    Members post updates on Messenger. When a problem or a risk of delay
    appears, they say so early so the team can adjust or help.

  \item \textbf{Manage shared work:}
    Documents are shared on Google Drive; source code is managed with
    Git and GitHub. Each member commits their own work with messages
    that describe the change clearly.

  \item \textbf{Cross-check:}
    When a member finishes a task, the other member reviews it for
    completeness, consistency and how well it meets the requirements
    before it is merged into the report.

  \item \textbf{Review before submission:}
    Both members check the product, the report, the structure of the
    submission and the commit evidence, so the submission is complete
    and on time.

  \item \textbf{Record actual contribution:}
    The task assignment and contribution percentages follow the work that
    was actually done, as agreed by both members.
\end{enumerate}

\noindent\textbf{Acceptance rule:}
A task is recorded as done only after both members have checked it and
agreed to accept it.
